% Section 5: Experimental methodology (+ Appendix B), converted from
% 05-experimental-methodology_4pages.docx
\documentclass[11pt]{article}
\usepackage[margin=2.5cm]{geometry}
\usepackage[T1]{fontenc}
\usepackage{lmodern}
\usepackage{amsmath,amssymb}
\usepackage{booktabs,tabularx,array}
\usepackage{caption}
\usepackage{chngcntr}
\usepackage[hidelinks]{hyperref}

\captionsetup{font=small,labelfont=bf}
\newcolumntype{L}{>{\raggedright\arraybackslash}X}
\newcolumntype{R}{>{\raggedleft\arraybackslash}X}
\counterwithin{table}{section}

\newcommand{\eps}{\varepsilon}
\newcommand{\Fhat}{\widehat{F}_0}
\newcommand{\dhat}{\widehat{\delta}}
\newcommand{\Err}{\mathsf{Error}}

\begin{document}

\setcounter{section}{4}
\section{Experimental methodology}\label{sec:methodology}

Section~2 proves that Algorithm~1 is an $(\eps,\delta)$-approximation with
bounded space and expected constant update time. The experiments test these
claims on streams with a known answer, measure how much of the proven margin
is used in practice, and compare the estimator with the baselines of
Section~3.2. This section fixes the design and the analysis; the results
follow in Section~6.

\subsection{Objectives}\label{sec:objectives}

The Project~9 brief asks for an accuracy study against $\eps$ and $\delta$, an
assessment of how conservative the threshold is, and a comparison with
classical estimators on synthetic and real streams. Each requirement becomes
one experiment (Table~\ref{tab:experiments}). E1 tests all three claims of
Theorem~1, accuracy, space and update time, on the same runs. Every run is
also checked against the analysis: the output is exact when $F_0 < \tau$, the
number of halvings matches the level predicted from $F_0$ and $\tau$, and the
estimate is unbiased (Table~\ref{tab:checks}). Table~\ref{tab:proof-steps}
states, for each experiment, the step of the proof it tests.

\begin{table}[htbp]
\centering\small
\caption{Experiments, the claims they test, and their status at the midterm
checkpoint.}\label{tab:experiments}
\begin{tabularx}{\linewidth}{@{}lLLl@{}}
\toprule
Exp. & Question & Claim tested & Status \\
\midrule
E1 & Does the estimator meet Theorem~1 in accuracy, space and update time? &
  $\Pr[\text{error} > \eps \text{ or } \bot] \le 5\delta/8$; at most
  $\tau = O(\eps^{-2}\log(m/\delta))$ stored items; expected amortised $O(1)$
  time per item (Sections 2.1--2.3) & run; $m$-sweep to run \\
E2 & How much smaller can $\tau$ be and still meet $\delta$? &
  slack in the Chernoff and union bounds (Section~2.3) & run \\
E3 & Which estimator is most accurate for equal memory, and which needs least
  memory for equal error? & comparison with WP4 baselines (Section~3.2) &
  awaits WP4 \\
E4 & Does stream order change the error? &
  guarantee holds for any fixed order (Section~1.1) & run (paper estimator) \\
E5 & Do the results hold on real streams? & E1 and E3 on real data &
  final report \\
\bottomrule
\end{tabularx}
\end{table}

\subsection{Estimators and workloads}\label{sec:estimators}

The \textbf{paper estimator} follows the corrected version of the paper
(Chakraborty, Vinodchandran and Meel, 2023), as implemented by WP2: each
arriving item is removed from the sample and re-inserted with probability $p$;
when the sample reaches $\tau$ elements, each is kept with probability $1/2$
and $p$ is halved; if none is removed, the run returns $\bot$. For E2 and E4
the threshold is scaled to $\tau_c = \max\{2, \lceil c\tau \rceil\}$; $c = 1$
is the algorithm of the paper. An \textbf{exact hash set} gives the true $F_0$
and is excluded from memory comparisons. The \textbf{baselines} of WP4
(HyperLogLog, $k$-minimum values and one off-the-shelf implementation) share
the same interface and a single seed. Memory is the logical state size:
$8\tau_c$ bytes for the paper estimator, the register array for HyperLogLog
and $k$ hash values for KMV. At a fixed budget of $B$ bytes the paper
estimator uses $\tau = B/8$, which no longer corresponds to a pair
$(\eps,\delta)$.

\begin{table}[htbp]
\centering\small
\caption{Workloads ($m = 100{,}000$, stream seed 0).}\label{tab:workloads}
\begin{tabularx}{\linewidth}{@{}lLrL@{}}
\toprule
Stream & Construction & $F_0$ & Purpose \\
\midrule
S1 uniform & $m$ draws from $\{1,\dots,10^5\}$ & 63,249 & moderate repetition \\
S2 all-distinct & random permutation of $1,\dots,m$ & 100,000 &
  largest $F_0$, most halvings \\
S3 Zipf & $m$ draws from $\{1,\dots,10^6\}$, $\Pr[k] \propto k^{-1.1}$ &
  23,999 & skewed, web-like; repeats must count once \\
S4 heavy repeats & 300 values repeated & 300 &
  $F_0 < \tau$: output must be exact (unit tests) \\
S5 reordered & S3 shuffled, sorted, in bursts, new items last & 23,999 &
  favourable and unfavourable orders (E4) \\
R1, R2 & two public logs & --- & real data (E5, final report) \\
\bottomrule
\end{tabularx}
\end{table}

We generate the streams ourselves with NumPy from a fixed seed, so each stream
and its exact $F_0$ are known and identical on every machine. Each stream is
fixed before any algorithm run, which is the setting of Theorem~1; an adaptive
adversary that sees the coin flips is not tested. Construction details are in
Table~\ref{tab:streams}.

\subsection{Parameters and runs}\label{sec:parameters}

The estimator only approximates when $F_0 \gg \tau$. For $m = 10^5$,
$\eps = 0.05$ would give $\tau > 110{,}000$, above every $F_0$, so we use
$\eps \in \{0.2, 0.3, 0.5\}$ and $\delta \in \{0.01, 0.1\}$, giving $\tau$
between 1,101 and 7,877 (6,880 for $\eps = 0.2, \delta = 0.1$) and two to
seven halvings. For the space and time claims, E1 also runs S2 at
$m \in \{10^4, 10^5, 10^6\}$ with $\eps = 0.2$, $\delta = 0.1$ and 200 seeds
each: $\tau$ then grows only with $\log m$ while $F_0$ grows with $m$. E2
halves $c$ from 1 to $1/256$, enough to locate where the failure rate crosses
$\delta$. E4 uses $c = 1$ and $c = 1/16$, because at $c = 1$ errors of about
1\% are too small to reveal an order effect. E3 compares the estimators at
equal memory budgets and at equal target errors, fixing one side of the
accuracy--memory trade-off at a time. Each configuration runs with 2,000
algorithm seeds: with no failures the 95\% upper bound on the failure
probability is about $3.8/n = 0.0019$, below both values of $\delta$.
Table~\ref{tab:setup} summarises the setup; the thresholds are listed in
Tables~\ref{tab:tau} and~\ref{tab:tauc}, and the 53 configurations (106,000
runs) in Table~\ref{tab:configs}.

\begin{table}[htbp]
\centering\small
\caption{Summary of the setup and the reason for each choice.}\label{tab:setup}
\begin{tabularx}{\linewidth}{@{}lLL@{}}
\toprule
Choice & Setting & Why \\
\midrule
Streams & S1--S3, $m = 10^5$ (S4 in tests, S5 for E4) &
  $F_0 \gg \tau$; runs take minutes \\
$\eps, \delta$ & $\{0.2, 0.3, 0.5\} \times \{0.01, 0.1\}$ &
  $\tau$ below every $F_0$ \\
Seeds & stream 0; algorithm 0--1,999 &
  Theorem~1 is over the coins for a fixed stream \\
Space, time & largest sample and time per item in every run;
  $m \in \{10^4, 10^5, 10^6\}$ on S2 & $\tau$ grows with $\log m$, $F_0$ with $m$ \\
Runs & 2,000 per configuration (E1: 36,000; E2: 54,000; E4: 16,000) &
  no failures gives $\delta^+ = 0.0019 < \delta$ \\
Failure & $\bot$ or relative error $> \eps$ & the event bounded by Theorem~1 \\
E2 & $c = 1, 1/2, \dots, 1/256$ at $\eps = 0.2, \delta = 0.1$ &
  locates where failures cross $\delta$ \\
E3 & equal memory; equal target error & one side of the trade-off fixed \\
E4 & $c = 1$ and $c = 1/16$ &
  at $c = 1$ errors are too small to show order effects \\
\bottomrule
\end{tabularx}
\end{table}

\subsection{Randomness and reproducibility}\label{sec:randomness}

The stream seed (fixed at 0) determines the data; the algorithm seed
determines the coin flips. Because only the algorithm seed varies, the spread
over 2,000 runs measures the algorithm's own randomness, which is exactly the
probability space of Theorem~1. Each run draws all random choices from
\texttt{random.Random(seed)}: a choice with probability $q$ succeeds if a
uniform $u \in [0,1)$ satisfies $u < q$. The sample is a Python
\texttt{dict}, so halvings do not depend on Python's randomised hashing, and a
run is a deterministic function of its configuration and seed. Every run is
appended as one row to a CSV file and never dropped, including runs that
return $\bot$; all tables and figures are generated from these files.
Implementation details are in Table~\ref{tab:impl}.

\subsection{Metrics and analysis}\label{sec:metrics}

For an output $\Fhat$ we record the relative error $|\Fhat - F_0|/F_0$, the
signed error (to detect bias), the largest sample size, the number of
halvings, the logical memory and the time per item. A run \emph{fails} if it
returns $\bot$ or its relative error exceeds $\eps$. The failure rate
$\dhat = k/n$ is reported with the upper end $\delta^+$ of its 95\% Wilson
interval, which remains meaningful when $k = 0$; a configuration is consistent
with the guarantee when $\delta^+ \le \delta$. Table~\ref{tab:metrics} defines
every metric and gives the Wilson bound.

\textbf{E1} reports, for accuracy, the mean, 95th-percentile and maximum error
relative to $\eps$, with histograms of the signed error. For space, it checks
that the largest sample never exceeds $\tau_c$, compares it with the $F_0$
items of the exact counter, and plots it against $\eps$, $\delta$ and $m$ next
to the predicted $\eps^{-2}\log(m/\delta)$ shape; space is counted in stored
64-bit identifiers, so the $\log$ of the universe size in the bit bound is a
constant. For time, it reports the median time per item against $\tau_c$
(using the E2 runs) and against $m$, and the amortised halving work
$h\tau_c/m$ coin flips per item, both of which should stay flat.
\textbf{E2} reports the conservatism factor $1/c^*$, where $c^*$ is the
smallest scale with $\delta^+ \le \delta$ on every stream. To locate the
slack, each run is classified by the events of Section~2.3 ($\bot$; too many
halvings, read from the recorded halvings; an error beyond $\eps$ otherwise),
and at $c = 1$ the frequency of each event over all runs is compared with its
bound (Table~\ref{tab:events}). The measured spread is also compared with the
binomial approximation $\sqrt{(1-p)/(pF_0)}$, which explains the numbers but
does not replace the proof. \textbf{E4} compares the error distributions of
the four orders; for the hash-based baselines the estimates must be identical,
which checks their implementation. \textbf{E3} plots error against memory and
reports the memory each estimator needs for a target error; update times are
compared only between implementations in the same language.

\subsection{Limitations}\label{sec:limitations}

The estimator is given the true stream length $m$, which many real streams do
not announce. Each workload is a single stream, so the guarantee is confirmed
on these streams rather than on all streams; the final version repeats E1 on
several stream seeds and at $m = 10^6$. The conservatism factor is measured
only at $\eps = 0.2, \delta = 0.1$. With 2,000 runs, failure probabilities
below about 0.002 cannot be distinguished from zero.

\subsection{Environment}\label{sec:environment}

The experiments ran with Python~3.13 on 2 CPU cores [TODO: machine model, CPU,
RAM, OS] using NumPy, pandas and Matplotlib; E1, E2 and E4 took about 12
minutes. The midterm runs used the WP3 pipeline at commit \texttt{d8722af},
whose estimator matches the shared WP2 implementation apart from input
validation. Configurations are frozen in \texttt{experiments/configs.py}, and
\texttt{make setup}, \texttt{make test}, \texttt{make experiments} and
\texttt{make analyze} reproduce all results from a clean checkout [TODO:
release tag].

% ---------------------------------------------------------------------------
\clearpage
\appendix
\setcounter{section}{1}
\section{Experimental details}\label{app:experimental-details}

\begin{table}[htbp]
\centering\small
\caption{The statement of the proof that each experiment tests.}\label{tab:proof-steps}
\begin{tabularx}{\linewidth}{@{}p{2.1cm}LL@{}}
\toprule
Exp. & Statement tested & What is checked \\
\midrule
E1 accuracy & $\Pr[\Err] \le 5\delta/8 < \delta$, where $\Err$ is returning
  $\bot$ or an estimate outside $[(1-\eps)F_0, (1+\eps)F_0]$ (Section~2.3) &
  the failure rate stays below $\delta$ for every stream and $(\eps,\delta)$;
  how much of $\eps$ is used \\
E1 space & the sample never holds more than
  $\tau = \lceil 12\eps^{-2}\log_2(8m/\delta) \rceil$ items (Section~2.1) &
  the bound holds in every run; the space does not grow with $F_0$; it scales
  as $\eps^{-2}$, $\log(1/\delta)$ and $\log m$ \\
E1 time & the expected amortised update time is $O(1)$ (Section~2.2) &
  the time per item and the halving work per item do not grow with $\tau_c$
  or $m$ \\
E2 & $\tau$ is set so that three worst-case bounds hold at once, each a
  Chernoff bound (Lemma~2) combined with a union bound over up to $m$
  positions or $m+1$ levels, the source of $\log_2(8m/\delta)$ &
  the smallest fraction $c$ of $\tau$ that still meets $\delta$, i.e.\ the
  slack these steps introduce \\
E3 & the estimator uses $O(\eps^{-2}\log(m/\delta))$ space &
  accuracy at equal memory and memory at equal target error, against the
  exact counter and the WP4 baselines; time per item \\
E4 & Theorem~1 holds for every stream fixed in advance, in any order and with
  repetitions (Section~1.1) & reordering the same items, including
  unfavourable orders, leaves the error distribution unchanged; for
  hash-based baselines the estimates are identical \\
E5 & as E1 and E3 & the conclusions hold on two real log streams (final
  report) \\
\bottomrule
\end{tabularx}
\end{table}

\begin{table}[htbp]
\centering\small
\caption{Construction of the synthetic streams (NumPy \texttt{default\_rng},
stream seed 0, $m = 100{,}000$).}\label{tab:streams}
\begin{tabularx}{\linewidth}{@{}p{2.8cm}L@{}}
\toprule
Stream & Details \\
\midrule
S1 uniform & Each draw is uniform on $\{1,\dots,10^5\}$. About $1/e$ of the
  values are never drawn, so $F_0 \approx (1 - 1/e)\cdot 10^5 \approx 63{,}212$;
  this stream has 63,249. \\
S2 all-distinct & A random permutation of $1,\dots,m$, so $F_0 = m$. \\
S3 Zipf & Draws from $\{1,\dots,10^6\}$ with $\Pr[k] \propto k^{-1.1}$; the
  probabilities are built explicitly because NumPy's Zipf sampler has no upper
  bound. Value 1 occurs 12,580 times; 18,943 values occur once. \\
S4 heavy repeats & 300 values repeated to length $m$, so $F_0 = 300 < \tau$.
  Used only in unit tests. \\
S5 reordered & The items of S3 in four orders: \emph{random} (shuffled),
  \emph{sorted} (ascending), \emph{bursts} (all copies of a value adjacent,
  values in random order), \emph{new items last} (most frequent values first,
  values occurring once at the end). $F_0$ is unchanged. \\
\bottomrule
\end{tabularx}
\end{table}

\begin{table}[htbp]
\centering\small
\begin{minipage}[t]{0.46\linewidth}
\centering
\caption{Thresholds $\tau$ from Equation~(2) for the E1 grid
($m = 100{,}000$).}\label{tab:tau}
\begin{tabular}{@{}lrr@{}}
\toprule
$\eps$ & $\delta = 0.01$ & $\delta = 0.1$ \\
\midrule
0.2 & 7,877 & 6,880 \\
0.3 & 3,501 & 3,058 \\
0.5 & 1,261 & 1,101 \\
\bottomrule
\end{tabular}
\end{minipage}\hfill
\begin{minipage}[t]{0.46\linewidth}
\centering
\caption{Scaled thresholds $\tau_c$ used in E2 ($\eps = 0.2$,
$\delta = 0.1$).}\label{tab:tauc}
\begin{tabular}{@{}lrlr@{}}
\toprule
$c$ & $\tau_c$ & $c$ & $\tau_c$ \\
\midrule
1    & 6,880 & 1/32  & 215 \\
1/2  & 3,440 & 1/64  & 108 \\
1/4  & 1,720 & 1/128 & 54 \\
1/8  & 860   & 1/256 & 27 \\
1/16 & 430   &       &    \\
\bottomrule
\end{tabular}
\end{minipage}
\end{table}

\begin{table}[htbp]
\centering\small
\caption{Experiment configurations; each runs with 2,000 algorithm
seeds.}\label{tab:configs}
\begin{tabularx}{\linewidth}{@{}lLrr@{}}
\toprule
Exp. & Parameters & Configs & Runs \\
\midrule
E1 & S1--S3; $\eps \in \{0.2, 0.3, 0.5\}$; $\delta \in \{0.01, 0.1\}$;
  $c = 1$ & 18 & 36,000 \\
E1 ($m$-sweep) & S2; $m \in \{10^4, 10^5, 10^6\}$; $\eps = 0.2$;
  $\delta = 0.1$; 200 seeds & 3 & 600 \\
E2 & S1--S3; $\eps = 0.2$; $\delta = 0.1$; $c \in \{2^0, \dots, 2^{-8}\}$ &
  27 & 54,000 \\
E4 & S5 in 4 orders; $\eps = 0.2$; $\delta = 0.1$; $c \in \{1, 1/16\}$ &
  8 & 16,000 \\
E3 & memory budgets and target errors fixed when the WP4 estimators are
  integrated & --- & --- \\
\midrule
Total (run) & & 53 & 106,000 \\
\bottomrule
\end{tabularx}
\end{table}

\begin{table}[htbp]
\centering\small
\caption{Metrics recorded for each run (output $\Fhat$, true count
$F_0$).}\label{tab:metrics}
\begin{tabularx}{\linewidth}{@{}p{3.4cm}L@{}}
\toprule
Metric & Definition \\
\midrule
Relative error & $|\Fhat - F_0|/F_0$, the quantity bounded by Equation~(1) \\
Signed error & $(\Fhat - F_0)/F_0$, used to detect bias \\
Guarantee failure & the run returns $\bot$ or its relative error exceeds
  $\eps$ \\
Failure rate & $\dhat = k/n$ over the $n$ runs of a configuration \\
Wilson upper bound &
  $\displaystyle \delta^+ = \frac{\dhat + \frac{z^2}{2n}
  + z\sqrt{\frac{\dhat(1-\dhat)}{n} + \frac{z^2}{4n^2}}}{1 + \frac{z^2}{n}}$,
  $z = 1.96$; about $3.8/n$ when $k = 0$ \\
Stored elements, halvings & largest sample size reached (at most $\tau_c$)
  and number of times $p$ was halved \\
Memory & logical memory under the common model (Section~\ref{sec:estimators}) \\
Halving work & $h\tau_c/m$, an upper bound on the coin flips per item spent
  in halvings \\
Update time & wall-clock time for the whole stream
  (\texttt{time.perf\_counter\_ns}) divided by $m$; stream generation and the
  exact count are not timed \\
\bottomrule
\end{tabularx}
\end{table}

\begin{table}[htbp]
\centering\small
\caption{Failure events of Section~2.3 and how each is identified in a
run.}\label{tab:events}
\begin{tabularx}{\linewidth}{@{}lLr@{}}
\toprule
Event & Identified by & Bound ($c = 1$) \\
\midrule
$\mathsf{Fail}$ & the run returns $\bot$ & $\delta/8$ \\
$\mathsf{Bad}_2$ & final $p = 2^{-h}$ below the cutoff, i.e.\ $h > \ell$ with
  $\ell = \lfloor \log_2(4F_0/\tau) \rfloor$ & $\delta/4$ \\
$\Err_2 \cap \neg\mathsf{Bad}_2$ & relative error $> \eps$ although
  $h \le \ell$ & $\delta/4$ \\
\bottomrule
\end{tabularx}
\end{table}

\begin{table}[htbp]
\centering\small
\caption{Consistency checks between each run and the analysis of Section~2;
a violation indicates an implementation error.}\label{tab:checks}
\begin{tabularx}{\linewidth}{@{}p{2.8cm}L@{}}
\toprule
Check & Expected behaviour \\
\midrule
Space & the largest sample size never exceeds $\tau_c$ (Section~2.1) \\
Exactness & when $F_0 < \tau$, no halving occurs and the output equals $F_0$
  (tested on S4) \\
Halvings & $h$ is close to $\lceil \log_2(F_0/\tau) \rceil$ and at most
  $\ell$; e.g.\ S2 with $\tau = 6{,}880$ gives 4 and $\ell = 5$ \\
Unbiasedness & the mean signed error is zero within its standard error
  (Appendix~A.1) \\
Update time & the median time per item does not grow with $\tau_c$ or with
  the number of halvings (Section~2.2) \\
\bottomrule
\end{tabularx}
\end{table}

\begin{table}[htbp]
\centering\small
\caption{Implementation and reproducibility.}\label{tab:impl}
\begin{tabularx}{\linewidth}{@{}p{2.8cm}L@{}}
\toprule
Item & Detail \\
\midrule
Random choices & each run creates \texttt{random.Random(seed)}; a choice with
  probability $q$ draws $u \in [0,1)$ and succeeds if $u < q$ \\
Sample storage & a Python \texttt{dict} (insertion order), so the elements
  removed in a halving do not depend on randomised hashing \\
Isolation & each run builds a fresh estimator; streams are cached per worker
  and never modified \\
Execution & runs are distributed over the CPU cores in chunks of 25 seeds \\
Storage & one CSV row per run (configuration, both seeds, true count,
  estimate, error, failure flag, stored elements, halvings, threshold, memory,
  time per item); raw files are never edited \\
Resumption & (configuration, seed) pairs already in the file are skipped, so
  an interrupted experiment restarts without duplicates \\
$\bot$ runs & recorded with the failure flag set and an empty estimate; never
  dropped \\
Commands & \texttt{make setup} (environment), \texttt{make test} (unit tests),
  \texttt{make experiments} (raw CSV files), \texttt{make analyze} (tables and
  figures) \\
\bottomrule
\end{tabularx}
\end{table}

\end{document}
